% ============================================================
%  GUÍA 10: Regla de la Cadena, Trascendentes e Implícita
%  Curso de Introducción a la Universidad — Precálculo y Cálculo I
%  Autor: Ing. Gabriel Astudillo
%  Clase cubierta: 10
% ============================================================

\documentclass[12pt, a4paper]{article}

% ── Paquetes ─────────────────────────────────────────────────
\usepackage{mathwizards-guia}
\mwguia{Guía 10 — Regla de la Cadena y otros}{Ing. Gabriel Astudillo $\cdot$ Curso Preuniversitario}

\begin{document}

% ── Portada / Título ─────────────────────────────────────────
\mwportada{Guía de Estudio 10}{Regla de la Cadena, Trascendentes e Implícita}{Regla de la Cadena, Derivadas de Exponenciales, Logarítmicas y Trigonométricas, Derivación Implícita y Orden Superior}

\vspace{4pt}
\begin{cuadroinfo}
  \small
  \textbf{Presentación:} Esta guía extiende la derivación a las funciones compuestas y trascendentes: la regla de la cadena, las derivadas de funciones exponenciales, logarítmicas y trigonométricas, la derivación implícita, las derivadas de orden superior y las funciones trigonométricas inversas. Con teoría, ejercicios resueltos y propuestos con clave.
  \\[4pt]
  \textbf{Nivel de Dificultad:}\quad
  \facil{} Básico / Directo \quad
  \medio{} Intermedio \quad
  \dificil{} Desafiante \quad
  \muyDificil{} Muy Difícil / Reto
\end{cuadroinfo}

\vspace{8pt}

\section{Regla de la Cadena}
% ============================================================

\subsection{¿Qué es una función compuesta?}

\begin{cuadroteoria}
  \textbf{Composición de funciones:} Si tenemos dos funciones $g$ y $f$, la \textbf{función compuesta} $(f \circ g)(x) = f(g(x))$ consiste en aplicar la función externa $f$ al resultado de la función interna $g$.

  Intuitivamente, es un proceso en dos etapas o una ``máquina de dos pasos'':
  \begin{center}
    \begin{tikzpicture}[
      box/.style={draw=mwprimario, fill=mwprimario!10, rounded corners=5pt, line width=1.1pt, minimum height=0.9cm, minimum width=2.4cm, align=center, font=\small\bfseries},
      arrow/.style={-{Stealth[scale=1.1]}, line width=1.1pt, draw=mwgrisoscuro}
      ]
      \node (in) at (0,0) {$x$};
      \node[box] (g) at (2.4,0) {Interior $g$\\[-1pt]\footnotesize (aplica $g$)};
      \node (mid) at (5.3,0) {$u = g(x)$};
      \node[box] (f) at (8.2,0) {Exterior $f$\\[-1pt]\footnotesize (aplica $f$)};
      \node (out) at (11.4,0) {$y = f(g(x))$};

      \draw[arrow] (in) -- (g) node[midway, above, font=\footnotesize\color{mwgris}] {Entrada};
      \draw[arrow] (g) -- (mid);
      \draw[arrow] (mid) -- (f);
      \draw[arrow] (f) -- (out) node[midway, above, font=\footnotesize\color{mwgris}] {Salida};
    \end{tikzpicture}
  \end{center}
\end{cuadroteoria}

\textbf{Identificación de capas (interna y externa):} Para derivar con éxito, el paso fundamental es reconocer la estructura $y = f(u)$ con $u = g(x)$:
\begin{center}
  \small
  \begin{tabular}{lll}
    \toprule
    \textbf{Función compuesta $y = f(g(x))$} & \textbf{Capa interna $u = g(x)$} & \textbf{Capa externa $y = f(u)$} \\
    \midrule
    $y = (3x^2 - 1)^5$                       & $u = 3x^2 - 1$                   & $y = u^5$                        \\
    $y = e^{2x}$                             & $u = 2x$                         & $y = e^u$                        \\
    $y = \ln(x^2 + 1)$                       & $u = x^2 + 1$                    & $y = \ln u$                      \\
    $y = \sin(x^3)$                          & $u = x^3$                        & $y = \sin u$                     \\
    $y = \sqrt{e^x + 1}$                     & $u = e^x + 1$                    & $y = \sqrt{u} = u^{1/2}$         \\
    \bottomrule
  \end{tabular}
\end{center}

\newpage

\subsection{Intuición de la Regla de la Cadena: Tasas de cambio engranadas}

¿Cómo cambia la salida $y$ cuando se altera la entrada $x$?

\begin{cuadroextra}
  \textbf{Analogía de las tasas de cambio (engranajes):}
  \begin{itemize}[leftmargin=*]
    \item Supongamos que $u$ cambia $2$ veces más rápido que $x$ ($\dfrac{\Delta u}{\Delta x} = 2$).
    \item A su vez, $y$ cambia $3$ veces más rápido que $u$ ($\dfrac{\Delta y}{\Delta u} = 3$).
    \item ¿A qué ritmo cambia $y$ respecto a $x$? Si $x$ avanza $1$ unidad, $u$ cambia en $2$ unidades, y esas $2$ unidades provocan que $y$ cambie en $2 \times 3 = 6$ unidades.
  \end{itemize}
  Las razones de cambio se \textbf{multiplican en cadena}:
  $$\frac{\Delta y}{\Delta x} = \frac{\Delta y}{\Delta u} \cdot \frac{\Delta u}{\Delta x}$$
  Al tomar el límite cuando $\Delta x \to 0$, obtenemos la célebre fórmula de Leibniz:
  $$\frac{dy}{dx} = \frac{dy}{du} \cdot \frac{du}{dx}$$
\end{cuadroextra}

\subsection{La Regla de la Cadena: Regla general y derivadas inmediatas}

\begin{cuadroteoria}
  \textbf{Regla de la cadena:} Si $y = f(g(x))$ y ambas funciones son derivables, entonces:
  $$\frac{dy}{dx} = f'(g(x)) \cdot g'(x)$$

  En notación de Leibniz, si $y = f(u)$ con $u = g(x)$:
  $$\frac{dy}{dx} = \frac{dy}{du} \cdot \frac{du}{dx}$$
\end{cuadroteoria}

\begin{cuadrosolucion}
  \textbf{Formas útiles de la regla de la cadena (con función interna $u(x)$):}
  \begin{align*}
    \frac{d}{dx}[(u(x))^n]   & = n(u(x))^{n-1} \cdot u'(x) & \quad \frac{d}{dx}[e^{u(x)}]   & = e^{u(x)} \cdot u'(x)     \\[4pt]
    \frac{d}{dx}[\ln(u(x))]  & = \frac{u'(x)}{u(x)}        & \quad \frac{d}{dx}[\sin(u(x))] & = \cos(u(x)) \cdot u'(x)   \\[4pt]
    \frac{d}{dx}[\cos(u(x))] & = -\sin(u(x)) \cdot u'(x)   & \quad \frac{d}{dx}[\tan(u(x))] & = \sec^2(u(x)) \cdot u'(x)
  \end{align*}
\end{cuadrosolucion}

\begin{cuadroadvertencia}
  \textbf{Estrategia y error común:}
  \begin{itemize}[leftmargin=*]
    \item \textbf{Estrategia:} Identifica la función de afuera $f$ y la de adentro $u$. Deriva la de afuera dejando la de adentro intacta, y multiplica al final por la derivada de adentro: $f'(u) \cdot u'$.
    \item \textbf{Error clásico:} Olvidar multiplicar por la derivada interna $u'(x)$ o derivar la de adentro antes de tiempo. Por ejemplo, la derivada de $\sin(x^3)$ \textbf{no} es $\cos(x^3)$, sino $\cos(x^3) \cdot 3x^2 = 3x^2 \cos(x^3)$.
  \end{itemize}
\end{cuadroadvertencia}

\subsection{Ejercicios resueltos}

\textbf{Ejemplo R1.} Calcula la derivada de $y = (3x^2 - 1)^5$.

\begin{cuadrosolucion}
  \textbf{Solución:} Función exterior: $u^5$ (derivada: $5u^4$). Función interior: $u = 3x^2 - 1$ (derivada: $6x$).
  $$y' = 5(3x^2 - 1)^4 \cdot 6x = 30x(3x^2 - 1)^4.$$
\end{cuadrosolucion}

\textbf{Ejemplo R2.} Calcula la derivada de $y = e^{2x}$.

\begin{cuadrosolucion}
  \textbf{Solución:} La función exterior es $e^u$ y la interior $u = 2x$:
  $$y' = e^{2x} \cdot 2 = 2e^{2x}.$$
\end{cuadrosolucion}

\textbf{Ejemplo R3.} Calcula la derivada de $y = \ln(x^2 + 1)$.

\begin{cuadrosolucion}
  \textbf{Solución:}
  $$y' = \frac{1}{x^2 + 1} \cdot 2x = \frac{2x}{x^2 + 1}.$$
\end{cuadrosolucion}

\textbf{Ejemplo R4.} Calcula la derivada de $y = \sin(x^3)$.

\begin{cuadrosolucion}
  \textbf{Solución:}
  $$y' = \cos(x^3) \cdot 3x^2 = 3x^2 \cos(x^3).$$
\end{cuadrosolucion}

\textbf{Ejemplo R5.} Calcula la derivada de $y = \sqrt{e^x + 1}$.

\begin{cuadrosolucion}
  \textbf{Solución:} Reescribimos $y = (e^x + 1)^{1/2}$:
  $$y' = \frac{1}{2}(e^x + 1)^{-1/2} \cdot e^x = \frac{e^x}{2\sqrt{e^x + 1}}.$$
\end{cuadrosolucion}

\newpage

\textbf{Ejemplo R6.} Calcula la derivada de $y = \cos^2(3x)$ (es decir, $[\cos(3x)]^2$).

\begin{cuadrosolucion}
  \textbf{Solución:} Aplicamos cadena dos veces (doble cadena):
  \begin{enumerate}
    \item Exterior: $u^2$ con $u = \cos(3x)$: $2\cos(3x)$.
    \item Interior: derivada de $\cos(3x) = -\sin(3x) \cdot 3 = -3\sin(3x)$.
  \end{enumerate}
  $$y' = 2\cos(3x) \cdot [-3\sin(3x)] = -6\cos(3x)\sin(3x) = -3\sin(6x).$$
  (En el último paso usamos $2\sin\theta\cos\theta = \sin(2\theta)$.)
\end{cuadrosolucion}


\subsection{Ejercicios propuestos}

\begin{enumerate}[series=ejercicios, leftmargin=*]
  \item \facil{} Calcula: $y'$ si $y = (2x + 1)^4$.
  \item \facil{} Calcula: $y'$ si $y = e^{3x}$.
  \item \facil{} Calcula: $y'$ si $y = \ln(x + 5)$.
  \item \facil{} Calcula: $y'$ si $y = \sin(4x)$.
  \item \medio{} Calcula: $y'$ si $y = (x^2 + 3x - 1)^6$.
  \item \medio{} Calcula: $y'$ si $y = e^{-x^2}$.
  \item \medio{} Calcula: $y'$ si $y = \ln(2x^3 + 1)$.
  \item \medio{} Calcula: $y'$ si $y = \cos(e^x)$.
  \item \dificil{} Calcula: $y'$ si $y = \sqrt{\sin x + 2}$.
  \item \dificil{} Calcula: $y'$ si $y = \tan(x^2 + 1)$.
  \item \dificil{} Calcula: $y'$ si $y = \left(\dfrac{e^x}{x}\right)^3$.
  \item \dificil{} Calcula: $y'$ si $y = \ln\left(\dfrac{x + 1}{x - 1}\right)$.
\end{enumerate}

\newpage

% ============================================================

\section{Derivadas de Orden Superior}
% ============================================================

\subsection{Notación e interpretación}

\begin{cuadroteoria}
  \textbf{Segunda derivada:} $f''(x) = \dfrac{d}{dx}[f'(x)] = \dfrac{d^2y}{dx^2}$

  \textbf{Tercera derivada:} $f'''(x) = \dfrac{d^3y}{dx^3}$

  \textbf{$n$-ésima derivada:} $f^{(n)}(x)$

  \textbf{Interpretación física:}
  \begin{itemize}[leftmargin=*]
    \item Si $s(t)$ es la posición, entonces $s'(t)$ es la velocidad y $s''(t)$ es la aceleración.
    \item Si $f''(x) > 0$ en un intervalo, $f$ es cóncava hacia arriba.
    \item Si $f''(x) < 0$ en un intervalo, $f$ es cóncava hacia abajo.
  \end{itemize}
\end{cuadroteoria}

\subsection{Ejercicios resueltos}

\textbf{Ejemplo R1.} Calcula la segunda derivada de $f(x) = x^3 - 4x^2 + 5$.

\begin{cuadrosolucion}
  \textbf{Solución:}
  \begin{align*}
    f'(x)  & = 3x^2 - 8x \\
    f''(x) & = 6x - 8.
  \end{align*}
\end{cuadrosolucion}

\textbf{Ejemplo R2.} Calcula la tercera derivada de $y = \sin x$.

\begin{cuadrosolucion}
  \textbf{Solución:}
  \begin{align*}
    y'   & = \cos x   \\
    y''  & = -\sin x  \\
    y''' & = -\cos x.
  \end{align*}
\end{cuadrosolucion}

\textbf{Ejemplo R3.} La posición de un objeto es $s(t) = t^3 - 6t^2 + 9t$ metros ($t$ en segundos). Encuentra velocidad y aceleración en $t = 2$.

\begin{cuadrosolucion}
  \textbf{Solución:}
  \begin{itemize}
    \item Velocidad: $v(t) = s'(t) = 3t^2 - 12t + 9$. En $t = 2$: $v(2) = 12 - 24 + 9 = -3$ m/s.
    \item Aceleración: $a(t) = v'(t) = s''(t) = 6t - 12$. En $t = 2$: $a(2) = 0$ m/s$^2$.
  \end{itemize}
  El objeto se mueve hacia atrás (velocidad negativa) pero la aceleración es cero en ese instante.
\end{cuadrosolucion}

\textbf{Ejemplo R4.} Calcula $f''(x)$ si $f(x) = x^2 e^x$.

\begin{cuadrosolucion}
  \textbf{Solución:}
  \begin{enumerate}
    \item Primera derivada (producto): $f' = 2x e^x + x^2 e^x = e^x(x^2 + 2x)$.
    \item Segunda derivada: derivamos $e^x(x^2 + 2x)$ con producto:
          $$f'' = e^x(x^2 + 2x) + e^x(2x + 2) = e^x(x^2 + 4x + 2).$$
  \end{enumerate}
\end{cuadrosolucion}

\subsection{Ejercicios propuestos}

\begin{enumerate}[resume=ejercicios, leftmargin=*]
  \item \facil{} Calcula $f''(x)$ para $f(x) = x^5 - 3x^2 + 1$.
  \item \facil{} Calcula $y''$ para $y = e^x$.
  \item \facil{} Calcula $y''$ para $y = \ln x$.
  \item \facil{} Calcula $y''$ para $y = \sin x$.
  \item \medio{} Calcula $f''(x)$ para $f(x) = x^3 \ln x$.
  \item \medio{} Calcula $y''$ para $y = e^{2x}$.
  \item \medio{} Calcula la tercera derivada de $f(x) = x^4 - 2x^3 + x^2$.
  \item \medio{} Si $s(t) = t^3 - 9t^2 + 24t$ (posición en metros, t en segundos), encuentra los instantes en que la velocidad es cero.
  \item \dificil{} Calcula $f'''(x)$ para $f(x) = \sin x \cdot \cos x$. (Pista: usa la identidad $2\sin x\cos x = \sin(2x)$.)
  \item \dificil{} Calcula la segunda derivada de $f(x) = \ln(x^2 + 1)$.
  \item \dificil{} Demuestra que $y = e^{2x}$ satisface la ecuación diferencial $y'' - 4y = 0$.
  \item \dificil{} Para $f(x) = \sqrt{x}$, calcula $f''(x)$ y determina los intervalos de concavidad.
\end{enumerate}

\newpage

% ============================================================

\section{Derivación Implícita}
% ============================================================

\subsection{Técnica}

\begin{cuadroteoria}
  \textbf{Derivación implícita:} Se usa cuando $y$ no está despejada explícitamente en función de $x$ (por ejemplo, $x^2 + y^2 = 25$).

  \textbf{Pasos:}
  \begin{enumerate}[leftmargin=*]
    \item Deriva ambos lados de la ecuación respecto a $x$.
    \item Cada vez que derives un término que contiene $y$, aplica la regla de la cadena: multiplica por $\dfrac{dy}{dx}$.
    \item Despeja $\dfrac{dy}{dx}$.
  \end{enumerate}
\end{cuadroteoria}

\begin{cuadroadvertencia}
  \textbf{Idea clave:} $y$ es una función de $x$, aunque no la conozcas explícitamente. Por tanto, derivar $y^2$ respecto a $x$ da $2y \cdot \frac{dy}{dx}$ (regla de la cadena), no $2y$.
\end{cuadroadvertencia}

\subsection{Ejercicios resueltos}

\textbf{Ejemplo R1.} Encuentra $\dfrac{dy}{dx}$ si $x^2 + y^2 = 25$.

\begin{cuadrosolucion}
  \textbf{Solución:}
  \begin{enumerate}
    \item Derivamos ambos lados: $2x + 2y\dfrac{dy}{dx} = 0$.
    \item Despejamos: $2y\dfrac{dy}{dx} = -2x \Rightarrow \dfrac{dy}{dx} = -\dfrac{x}{y}$.
  \end{enumerate}
  Nota: la pendiente depende del punto $(x, y)$ en la circunferencia.
\end{cuadrosolucion}

\newpage

\textbf{Ejemplo R2.} Encuentra $\dfrac{dy}{dx}$ si $x^3 + y^3 = 6xy$.

\begin{cuadrosolucion}
  \textbf{Solución:}
  \begin{enumerate}
    \item Derivamos: $3x^2 + 3y^2 \dfrac{dy}{dx} = 6y + 6x\dfrac{dy}{dx}$.
    \item Agrupamos los términos con $\dfrac{dy}{dx}$: $3y^2 \dfrac{dy}{dx} - 6x\dfrac{dy}{dx} = 6y - 3x^2$.
    \item Factorizamos: $\dfrac{dy}{dx}(3y^2 - 6x) = 6y - 3x^2$.
    \item $\dfrac{dy}{dx} = \dfrac{6y - 3x^2}{3y^2 - 6x} = \dfrac{2y - x^2}{y^2 - 2x}$.
  \end{enumerate}
\end{cuadrosolucion}

\textbf{Ejemplo R3.} Encuentra la ecuación de la recta tangente a $x^2 + y^2 = 25$ en el punto $(3, 4)$.

\begin{cuadrosolucion}
  \textbf{Solución:}
  \begin{enumerate}
    \item Del Ejemplo R1: $\dfrac{dy}{dx} = -\dfrac{x}{y}$.
    \item En $(3, 4)$: $m = -\dfrac{3}{4}$.
    \item Recta tangente: $y - 4 = -\dfrac{3}{4}(x - 3) \Rightarrow y = -\dfrac{3}{4}x + \dfrac{25}{4}$.
  \end{enumerate}
\end{cuadrosolucion}

\textbf{Ejemplo R4.} Encuentra $\dfrac{dy}{dx}$ si $\sin(xy) = x + y$.

\begin{cuadrosolucion}
  \textbf{Solución:}
  \begin{enumerate}
    \item Derivamos (con cadena en $\sin(xy)$ y producto en $xy$):
          $$\cos(xy) \cdot \left(y + x\dfrac{dy}{dx}\right) = 1 + \dfrac{dy}{dx}.$$
    \item Expandimos: $y\cos(xy) + x\cos(xy)\dfrac{dy}{dx} = 1 + \dfrac{dy}{dx}$.
    \item Agrupamos: $\dfrac{dy}{dx}\left[x\cos(xy) - 1\right] = 1 - y\cos(xy)$.
    \item $\dfrac{dy}{dx} = \dfrac{1 - y\cos(xy)}{x\cos(xy) - 1}$.
  \end{enumerate}
\end{cuadrosolucion}

\newpage

\subsection{Ejercicios propuestos}

\begin{enumerate}[resume=ejercicios, leftmargin=*]
  \item \facil{} Encuentra $\dfrac{dy}{dx}$ si $x^2 + y^2 = 16$.
  \item \facil{} Encuentra $\dfrac{dy}{dx}$ si $3x + 2y = 10$.
  \item \facil{} Encuentra $\dfrac{dy}{dx}$ si $xy = 12$.
  \item \medio{} Encuentra $\dfrac{dy}{dx}$ si $x^2 y + y^3 = 8$.
  \item \medio{} Encuentra $\dfrac{dy}{dx}$ si $x^2 + xy + y^2 = 7$.
  \item \medio{} Encuentra la ecuación de la recta tangente a $x^2 + y^2 = 25$ en $(4, -3)$.
  \item \medio{} Encuentra $\dfrac{dy}{dx}$ si $e^y = x + y$.
  \item \medio{} Encuentra $\dfrac{dy}{dx}$ si $\ln y + x = y^2$.
  \item \dificil{} Encuentra $\dfrac{dy}{dx}$ si $\cos(x + y) = xy$.
  \item \dificil{} Encuentra $\dfrac{dy}{dx}$ si $x^2 y^3 - 2x = y$.
  \item \dificil{} Encuentra $\dfrac{d^2y}{dx^2}$ (segunda derivada implícita) si $x^2 + y^2 = 25$.
  \item \dificil{} Una elipse tiene ecuación $\dfrac{x^2}{9} + \dfrac{y^2}{4} = 1$. Encuentra la pendiente de la recta tangente en el punto $\left(\dfrac{3}{\sqrt{2}},\; \sqrt{2}\right)$.
\end{enumerate}

\newpage

% ============================================================

\section{Derivadas de Funciones Trigonométricas Inversas}
% ============================================================

\subsection{Deducción y fórmulas}

\begin{cuadroteoria}
  \textbf{Derivada de $\arcsin x$} (deducción por derivación implícita):

  Sea $y = \arcsin x$, es decir, $\sin y = x$ con $-\frac{\pi}{2} \leq y \leq \frac{\pi}{2}$.
  \begin{enumerate}[leftmargin=*]
    \item Derivamos implícitamente: $\cos y \cdot \dfrac{dy}{dx} = 1$.
    \item $\dfrac{dy}{dx} = \dfrac{1}{\cos y}$.
    \item Como $\sin^2 y + \cos^2 y = 1$, entonces $\cos y = \sqrt{1 - \sin^2 y} = \sqrt{1 - x^2}$ (positivo porque $y \in [-\frac{\pi}{2}, \frac{\pi}{2}]$).
  \end{enumerate}
  $$\boxed{\frac{d}{dx}[\arcsin x] = \frac{1}{\sqrt{1 - x^2}}}$$

  \textbf{Derivada de $\arctan x$} (deducción similar):

  Sea $y = \arctan x$, es decir, $\tan y = x$.
  \begin{enumerate}[leftmargin=*]
    \item Derivamos: $\sec^2 y \cdot \dfrac{dy}{dx} = 1$.
    \item $\dfrac{dy}{dx} = \dfrac{1}{\sec^2 y} = \dfrac{1}{1 + \tan^2 y} = \dfrac{1}{1 + x^2}$.
  \end{enumerate}
  $$\boxed{\frac{d}{dx}[\arctan x] = \frac{1}{1 + x^2}}$$
\end{cuadroteoria}

\begin{cuadrosolucion}
  \textbf{Tabla de derivadas de inversas trigonométricas:}
  \begin{center}
    \begin{tabular}{ll}
      \toprule
      \textbf{Función} & \textbf{Derivada}            \\
      \midrule
      $\arcsin x$      & $\dfrac{1}{\sqrt{1 - x^2}}$  \\
      $\arccos x$      & $-\dfrac{1}{\sqrt{1 - x^2}}$ \\
      $\arctan x$      & $\dfrac{1}{1 + x^2}$         \\
      \bottomrule
    \end{tabular}
  \end{center}
  Nota: $\dfrac{d}{dx}[\arccos x] = -\dfrac{d}{dx}[\arcsin x]$ porque $\arcsin x + \arccos x = \dfrac{\pi}{2}$.
\end{cuadrosolucion}

\subsection{Ejercicios resueltos}

\textbf{Ejemplo R1.} Calcula la derivada de $y = \arcsin(3x)$.

\begin{cuadrosolucion}
  \textbf{Solución:} Aplicamos cadena con $u = 3x$:
  $$y' = \frac{1}{\sqrt{1 - (3x)^2}} \cdot 3 = \frac{3}{\sqrt{1 - 9x^2}}.$$
\end{cuadrosolucion}

\textbf{Ejemplo R2.} Calcula la derivada de $y = \arctan(x^2)$.

\begin{cuadrosolucion}
  \textbf{Solución:} Con $u = x^2$:
  $$y' = \frac{1}{1 + (x^2)^2} \cdot 2x = \frac{2x}{1 + x^4}.$$
\end{cuadrosolucion}

\textbf{Ejemplo R3.} Calcula la derivada de $y = x \arctan x$.

\begin{cuadrosolucion}
  \textbf{Solución:} Regla del producto:
  $$y' = \arctan x + x \cdot \frac{1}{1 + x^2} = \arctan x + \frac{x}{1 + x^2}.$$
\end{cuadrosolucion}

%\newpage

\subsection{Ejercicios propuestos}

\begin{enumerate}[resume=ejercicios, leftmargin=*]
  \item \facil{} Calcula: $y'$ si $y = \arcsin(2x)$.
  \item \facil{} Calcula: $y'$ si $y = \arctan(5x)$.
  \item \facil{} Calcula: $y'$ si $y = \arccos(x)$.
  \item \medio{} Calcula: $y'$ si $y = \arcsin(x^2)$.
  \item \medio{} Calcula: $y'$ si $y = \arctan\left(\dfrac{1}{x}\right)$.
  \item \medio{} Calcula: $y'$ si $y = \ln(\arctan x)$.
  \item \medio{} Calcula: $y'$ si $y = e^{\arcsin x}$.
  \item \dificil{} Calcula: $y'$ si $y = \arctan(\ln x)$.
  \item \dificil{} Calcula: $y'$ si $y = x^2 \arcsin x + \sqrt{1 - x^2}$. (Pista: la simplificación es elegante.)
  \item \dificil{} Encuentra la ecuación de la recta tangente a $y = \arctan x$ en $x = 1$.
\end{enumerate}

\newpage

% ============================================================

% ──────────────────────────────────────────────────────────
%  NUEVOS EJERCICIOS PROPUESTOS (4) — INSERCIÓN MANUAL
% ──────────────────────────────────────────────────────────
\subsection{Ejercicios propuestos: Retos Integradores}

\begin{enumerate}[resume=ejercicios, leftmargin=*]
  \item \medio{} Deriva aplicando la regla de la cadena: $f(x) = e^{x^2}$.

  \item \dificil{} Deriva la función compuesta: $f(x) = \ln(\sin x)$.

  \item \dificil{} Halla $y'$ por derivación implícita: $x^3 + y^3 = 6xy$.

  \item \muyDificil{} Halla la segunda derivada de $f(x) = x^3 \ln x$.
\end{enumerate}
\newpage
% ============================================================
\ifmwclaves
\section{Clave de Respuestas}
% ============================================================

\subsection*{Sección 1: Regla de la Cadena (Ejercicios 1--12)}

\begin{enumerate}[series=respuestas, leftmargin=*, label=\textbf{\arabic*.}]
  \item $y' = 4(2x + 1)^3 \cdot 2 = 8(2x + 1)^3$
  \item $y' = 3e^{3x}$
  \item $y' = \dfrac{1}{x + 5}$
  \item $y' = 4\cos(4x)$
  \item $y' = 6(x^2 + 3x - 1)^5(2x + 3)$
  \item $y' = e^{-x^2}(-2x) = -2x e^{-x^2}$
  \item $y' = \dfrac{6x^2}{2x^3 + 1}$
  \item $y' = -\sin(e^x) \cdot e^x = -e^x \sin(e^x)$
  \item $y' = \dfrac{\cos x}{2\sqrt{\sin x + 2}}$
  \item $y' = \sec^2(x^2 + 1) \cdot 2x = 2x \sec^2(x^2 + 1)$
  \item $y' = 3\left(\dfrac{e^x}{x}\right)^2 \cdot \left(\dfrac{e^x x - e^x}{x^2}\right) = 3\left(\dfrac{e^x}{x}\right)^2 \cdot \dfrac{e^x(x - 1)}{x^2} = \dfrac{3e^{3x}(x - 1)}{x^4}$
  \item $y = \ln(x + 1) - \ln(x - 1) \Rightarrow y' = \dfrac{1}{x + 1} - \dfrac{1}{x - 1} = \dfrac{(x - 1) - (x + 1)}{(x + 1)(x - 1)} = \dfrac{-2}{x^2 - 1}$
\end{enumerate}


\subsection*{Sección 2: Derivadas de Orden Superior (Ejercicios 13--24)}

\begin{enumerate}[resume=respuestas, leftmargin=*, label=\textbf{\arabic*.}]
  \item $f'(x) = 5x^4 - 6x$; $f''(x) = 20x^3 - 6$.
  \item $y' = e^x$; $y'' = e^x$.
  \item $y' = \dfrac{1}{x}$; $y'' = -\dfrac{1}{x^2}$.
  \item $y' = \cos x$; $y'' = -\sin x$.
  \item $f' = 3x^2 \ln x + x^3 \cdot \frac{1}{x} = 3x^2 \ln x + x^2$. $f'' = 6x \ln x + 3x^2 \cdot \frac{1}{x} + 2x = 6x \ln x + 3x + 2x = 6x\ln x + 5x$.
  \item $y' = 2e^{2x}$; $y'' = 4e^{2x}$.
  \item $f' = 4x^3 - 6x^2 + 2x$; $f'' = 12x^2 - 12x + 2$; $f''' = 24x - 12$.
  \item $v(t) = 3t^2 - 18t + 24 = 3(t^2 - 6t + 8) = 3(t - 2)(t - 4) = 0$. Velocidad cero en $t = 2$ y $t = 4$ (segundos).
  \item Usando la identidad: $f(x) = \frac{1}{2}\sin(2x)$. Entonces $f'(x) = \cos(2x)$, $f''(x) = -2\sin(2x)$, $f'''(x) = -4\cos(2x)$.
  \item $f'(x) = \dfrac{2x}{x^2 + 1}$; $f''(x) = \dfrac{2(x^2 + 1) - 2x(2x)}{(x^2 + 1)^2} = \dfrac{2 - 2x^2}{(x^2 + 1)^2}$.
  \item $y = e^{2x} \Rightarrow y' = 2e^{2x} \Rightarrow y'' = 4e^{2x}$. Entonces $y'' - 4y = 4e^{2x} - 4e^{2x} = 0$. Se verifica.
  \item $f(x) = x^{1/2}$; $f'(x) = \frac{1}{2}x^{-1/2}$; $f''(x) = -\frac{1}{4}x^{-3/2} = -\dfrac{1}{4\sqrt{x^3}} < 0$ para todo $x > 0$. La función es cóncava hacia abajo en $(0, \infty)$.
\end{enumerate}

\newpage


\subsection*{Sección 3: Derivación Implícita (Ejercicios 25--36)}

\begin{enumerate}[resume=respuestas, leftmargin=*, label=\textbf{\arabic*.}]
  \item $\dfrac{dy}{dx} = -\dfrac{x}{y}$
  \item $\dfrac{dy}{dx} = -\dfrac{3}{2}$
  \item $y + x\dfrac{dy}{dx} = 0 \Rightarrow \dfrac{dy}{dx} = -\dfrac{y}{x}$
  \item $2xy + x^2\dfrac{dy}{dx} + 3y^2\dfrac{dy}{dx} = 0 \Rightarrow \dfrac{dy}{dx} = -\dfrac{2xy}{x^2 + 3y^2}$
  \item $2x + y + x\dfrac{dy}{dx} + 2y\dfrac{dy}{dx} = 0 \Rightarrow \dfrac{dy}{dx} = -\dfrac{2x + y}{x + 2y}$
  \item $\dfrac{dy}{dx} = -\dfrac{x}{y}$; en $(4, -3)$: $m = \dfrac{4}{3}$. Recta: $y + 3 = \dfrac{4}{3}(x - 4) \Rightarrow y = \dfrac{4}{3}x - \dfrac{25}{3}$.
  \item $e^y \dfrac{dy}{dx} = 1 + \dfrac{dy}{dx} \Rightarrow \dfrac{dy}{dx}(e^y - 1) = 1 \Rightarrow \dfrac{dy}{dx} = \dfrac{1}{e^y - 1}$
  \item $\dfrac{1}{y}\dfrac{dy}{dx} + 1 = 2y\dfrac{dy}{dx} \Rightarrow \dfrac{dy}{dx}\left(\dfrac{1}{y} - 2y\right) = -1 \Rightarrow \dfrac{dy}{dx} = \dfrac{y}{2y^2 - 1}$
  \item $-\sin(x+y)\left(1 + \dfrac{dy}{dx}\right) = y + x\dfrac{dy}{dx} \Rightarrow \dfrac{dy}{dx} = \dfrac{-y - \sin(x+y)}{x + \sin(x+y)}$
  \item $2xy^3 + 3x^2 y^2 \dfrac{dy}{dx} - 2 = \dfrac{dy}{dx} \Rightarrow \dfrac{dy}{dx}(3x^2 y^2 - 1) = 2 - 2xy^3 \Rightarrow \dfrac{dy}{dx} = \dfrac{2 - 2xy^3}{3x^2 y^2 - 1}$
  \item $\dfrac{dy}{dx} = -\dfrac{x}{y}$. Derivamos de nuevo: $\dfrac{d^2y}{dx^2} = -\dfrac{y - x\frac{dy}{dx}}{y^2} = -\dfrac{y + \frac{x^2}{y}}{y^2} = -\dfrac{y^2 + x^2}{y^3} = -\dfrac{25}{y^3}$.
  \item $\dfrac{2x}{9} + \dfrac{2y}{4}\dfrac{dy}{dx} = 0 \Rightarrow \dfrac{dy}{dx} = -\dfrac{4x}{9y}$. En $\left(\frac{3}{\sqrt{2}}, \sqrt{2}\right)$: $m = -\dfrac{4 \cdot \frac{3}{\sqrt{2}}}{9\sqrt{2}} = -\dfrac{\frac{12}{\sqrt{2}}}{9\sqrt{2}} = -\dfrac{12}{18} = -\dfrac{2}{3}$.
\end{enumerate}


\subsection*{Sección 4: Inversas Trigonométricas (Ejercicios 37--46)}

\begin{enumerate}[resume=respuestas, leftmargin=*, label=\textbf{\arabic*.}]
  \item $y' = \dfrac{2}{\sqrt{1 - 4x^2}}$
  \item $y' = \dfrac{5}{1 + 25x^2}$
  \item $y' = -\dfrac{1}{\sqrt{1 - x^2}}$
  \item $y' = \dfrac{2x}{\sqrt{1 - x^4}}$
  \item $y' = \dfrac{1}{1 + (1/x)^2} \cdot \left(-\dfrac{1}{x^2}\right) = \dfrac{-1/x^2}{1 + 1/x^2} = -\dfrac{1}{x^2 + 1}$
  \item $y' = \dfrac{1}{\arctan x} \cdot \dfrac{1}{1 + x^2} = \dfrac{1}{(1 + x^2)\arctan x}$
  \item $y' = e^{\arcsin x} \cdot \dfrac{1}{\sqrt{1 - x^2}} = \dfrac{e^{\arcsin x}}{\sqrt{1 - x^2}}$
  \item $y' = \dfrac{1}{1 + (\ln x)^2} \cdot \dfrac{1}{x} = \dfrac{1}{x(1 + \ln^2 x)}$
  \item $y' = 2x\arcsin x + \dfrac{x^2}{\sqrt{1 - x^2}} + \dfrac{-x}{\sqrt{1-x^2}} = 2x\arcsin x + \dfrac{x^2 - x}{\sqrt{1 - x^2}} = 2x \arcsin x + \dfrac{x(x-1)}{\sqrt{1-x^2}}$.
  \item En $x = 1$: $y(1) = \arctan 1 = \frac{\pi}{4}$; $y'(1) = \frac{1}{1+1} = \frac{1}{2}$. Recta: $y - \frac{\pi}{4} = \frac{1}{2}(x - 1)$.
\end{enumerate}


% ──────────────────────────────────────────────────────────
%  NUEVAS RESPUESTAS (4) — INSERCIÓN MANUAL
% ──────────────────────────────────────────────────────────
\subsection*{Sección 5: Retos Integradores (Ejercicios 47--50)}
\begin{enumerate}[resume=respuestas, leftmargin=*, label=\textbf{\arabic*.}]
  \item $f'(x) = 2x\,e^{x^2}$
  \item $f'(x) = \dfrac{\cos x}{\sin x} = \cot x$
  \item $3x^2 + 3y^2 y' = 6y + 6x y' \Rightarrow y' = \dfrac{2y - x^2}{y^2 - 2x}$
  \item $f'(x) = 3x^2 \ln x + x^2$; \quad $f''(x) = 6x \ln x + 5x$
\end{enumerate}
\fi
\end{document}
